\ifdefined\MyWebBook\else
\documentclass[../textbook.tex]{subfiles}
\begin{document}
\fi
\bookchapter{绪论}{ch:introduction}

语言模型依据已有上下文预测后续符号的条件概率分布。通过设计输入输出表示，自回归生成可推广至问答、翻译、代码编写与工具调用。构建完整的模型系统需要协同语料准备、参数训练、推理解码策略、硬件资源规划与质量评估。

语言建模的演进交织着表示方法、网络结构与训练目标的共同演变。规模定律刻画了计算资源投入与验证损失之间的经验关系，基准评测考察模型在特定任务上的表现，而模型架构、优化算法与应用系统则分别承担概率映射、参数学习与任务执行的职责。

\section{语言建模的研究对象}

\subsection{文本序列化}
文本由字符组成，模型通常处理词元（Token）序列。词元是编码规则规定的离散单元，可以对应字节、字符、词的一部分或特定控制标记。给定有限词表 $\mathcal V$，一段文本经分词器转换为
\begin{equation}
 x_{1:T}=(x_1,x_2,\ldots,x_T),\qquad x_t\in\mathcal V.
\end{equation}
$T$ 是编码后的序列长度，不必等于字符数或词数。词元编号用于定位词表条目，词元间的语义关系由模型学习到的连续向量表示描述。

例如，“学习语言模型”可以被某种分词规则切分为“学习”“语言”“模型”三个词元，也可以按子词或单字切分成更多单元。不同的分词方案均能无损重构原文，但词元序列长度、词表大小与每步预测空间显著不同。因此，评估上下文窗口有效利用率、困惑度及推理吞吐时，必须基于确定的分词规则与词表定义。

语言模型给序列赋予概率。设参数集合为 $\theta$，模型分布记为 $p_\theta$。概率链式法则允许将联合概率分解为
\begin{equation}\label{eq:intro-chain}
 p_\theta(x_{1:T})=\prod_{t=1}^{T}p_\theta(x_t\mid x_{<t}),
 \qquad x_{<t}=(x_1,\ldots,x_{t-1}).
\end{equation}
第一个位置的条件前缀为空；实际系统也可以用约定的起始标记表示初始条件。式中每一项代表在给定历史前缀条件下的下一个词元概率分布，在词表空间上满足非负性与归一化约束。

式\eqref{eq:intro-chain}本身是概率乘法法则的恒等变换。神经网络参数化条件概率、受限上下文窗口与跨步参数共享则属于具体的模型假设；同一分解形式可由前馈网络、循环网络或 Transformer 等不同结构实现。

\subsection{任务的条件生成表示}
在问答或翻译中，输入条件 $c$ 与待生成序列 $y_{1:T}$ 具有不同职责。模型描述的是
\begin{equation}
 p_\theta(y_{1:T}\mid c)=\prod_{t=1}^{T}p_\theta(y_t\mid c,y_{<t}).
\end{equation}
$c$ 可以包含提示文本、对话历史、检索文档或多模态特征。条件表示的形式决定了输入信息的编码通路：离散提示文本经过嵌入查表，连续多模态特征则通过投影层对齐到语义空间。训练目标进一步约束预测分布：即使网络相同，仅在回答片段计算损失的监督微调，与全序列无差别自回归训练会产生截然不同的生成行为。

生成过程需要明确的终止条件。模型输出预定义的结束词元（End of Sequence，EOS）表示自主完成回答；运行时设置的最大步长则构成外部安全截断。每步概率归一化只约束单个位置的分布，当模型陷入重复循环或概率漂移时可能无法触发结束词元，系统须配备超时截断机制。

\begin{figure}[htbp]
\centering
\includegraphics[width=\linewidth]{\subfix{../../figures/theory/intro-autoregressive-loop.pdf}}
\caption[自回归生成与条件概率闭环]{自回归生成与条件概率闭环。模型基于当前前缀计算下一词元条件概率分布，经采样决策生成新词元；非终止词元拼接入前缀触发下一步计算，终止标记或达到最大长度时退出循环。}\label{fig:intro-autoregressive-loop}
\end{figure}


\section{语言模型的发展历程}

\label{sec:intro-history}
语言模型的发展贯穿三条线索：怎样表示上下文，怎样利用语料学习，以及怎样把预测能力用于任务。这三者长期并存，并在表示学习、预训练和任务适配上彼此推动。

\begin{figure}[p]
\centering
\includegraphics[width=\linewidth,height=0.82\textheight,keepaspectratio]{\subfix{../../figures/theory/intro-model-technology-timeline.pdf}}
\caption[基础模型与关键技术演进节点]{基础模型与关键技术演进节点。左侧为方法与架构，右侧为基础模型；方括号表示首次公开时的主要发布形态，紫色空心节点表示新近进展。}\label{fig:intro-model-technology-timeline}
\end{figure}

图\ref{fig:intro-model-technology-timeline}将基础模型家族与产生能力的技术路径解耦：同一模型可组合多种算法机制，同一机制亦能被不同模型复用。图谱限定于基础模型及其核心算法，应用产品与智能体运行时属于上层系统架构。节点依据公开文献与官方技术资料整理\cite{vaswani2017attention,radfordImprovingLanguageUnderstanding2018,devlinBERTPretrainingDeep2019,brown2020fewshot,kaplan2020scaling,ouyang2022instructgpt,hoffmann2022chinchilla,deepseek2024v2,deepseek2024v3,deepseek2025r1,anthropic2026fable51launch,deepseek2026v41flash}。

\subsection{分布式词表示}
Shannon 在1948年的通信理论工作中讨论了语言序列的统计近似，为将离散符号串纳入概率建模框架奠定了理论基础\cite{shannon1948communication}。早期统计语言模型常采用\term{$n$元语言模型}{$n$-gram Language Model}{}，基于马尔可夫假设用最近的 $n-1$ 个词近似完整历史前缀：
\begin{equation}
 p(x_t\mid x_{<t})\approx p(x_t\mid x_{t-n+1:t-1}).
\end{equation}
设已为序列开头填充约定的起始标记。给定历史上下文 $h$，可通过经验计数比 $\frac{c(h,w)}{c(h)}$ 估计后继词为 $w$ 的条件概率。当历史未曾在语料中出现时分母为零；当特定接续未见时直接计数会导致零概率分配。平滑、回退与插值算法通过在不同阶数上下文间分配概率质量以缓解零频问题。

增大 $n$ 虽能捕获更长距离的局部依赖，但参数组合空间随上下文长度指数级膨胀。有限语料难以覆盖长程组合，且离散符号统计无法泛化至语义相近的近义表达。Bengio 等人在2003年提出神经概率语言模型，联合学习连续词向量与条件概率分布\cite{bengio2003neuralprobabilistic}。模型通过可学习的连续向量表征共享上下文信息，有效缓解了传统 $n$ 元计数在高阶组合下的稀疏性问题，但仍受限于固定宽度的前馈输入窗口。

\subsection{循环网络的历史表示}
\term{循环神经网络}{Recurrent Neural Network}{RNN}通过递归状态方程将变长历史编码至固定维度隐向量中。Mikolov 等人在2010年提出的循环语言模型拓展了变长依赖建模能力\cite{mikolov2010rnnlm}。然而早期上下文信息需经过多步非线性状态转换才能作用于远端预测，反向传播过程中连续矩阵连乘容易引发梯度衰减或梯度爆炸。

\term{长短期记忆}{Long Short-Term Memory}{LSTM}网络通过引入门控机制与加性记忆单元，改善了跨时间步的信息保持与梯度反向传播\cite{hochreiter1997lstm}。门控单元有效延缓了长程记忆衰减，但固定容量的隐状态依然构成信息压缩瓶颈，且严格的时间步递归更新阻碍了序列维度的训练并行化。为消除串行依赖并支持任意位置间的直接交互，注意力机制成为后续演进的核心方向。

\subsection{预训练模型的发展}
利用大规模无标注语料预训练通用表示，再迁移至下游任务，成为现代自然语言处理的标准范式。静态词向量为每个词型分配固定嵌入，无法随上下文消解一词多义。2018年提出的 ELMo 利用双向长短期记忆语言模型的内部层激活拼接生成动态上下文表示\cite{peters2018elmo}。

Vaswani 等人在2017年提出 Transformer 架构，彻底摒弃循环连接，完全依赖自注意力机制计算序列内任意两个位置的关联分数，实现训练阶段全序列位置的高度并行计算\cite{vaswani2017attention}。这种并行性仅存在于\emph{训练阶段}；全注意力矩阵在长序列下具有 $\mathcal{O}(T^2)$ 的时间和空间复杂度。

2018 年的 GPT-1 采用因果单向自回归进行生成式预训练，再针对下游任务微调\cite{radfordImprovingLanguageUnderstanding2018}；BERT 则采用双向掩码语言建模（Masked Language Modeling，MLM）学习全局上下文表征\cite{devlinBERTPretrainingDeep2019}。二者骨干网络虽同为 Transformer，但其自注意力掩码可见性与概率建模目标截然不同：BERT 重建被遮蔽词元，GPT 建模因果序列生成。

\subsection{大规模语言模型的发展}
Brown 等人在2020年通过 GPT-3 展示了无需调整模型参数的\term{上下文学习}{In-Context Learning}{ICL}能力\cite{brown2020fewshot}。模型通过提示词中的任务说明与少样本示范直接完成目标任务。上下文学习改变的是前向推理时的\emph{条件信息}，而非\emph{模型权重}。

规模化规律进一步引导算力资源分配：在固定计算预算下，应协同扩大模型参数量与训练数据量。Kaplan 等人与 Hoffmann 等人（Chinchilla）通过对数多项式经验拟合，系统建立了计算最优配比关系\cite{kaplan2020scaling,hoffmann2022chinchilla}。最佳配比参数取决于硬件通信效率、语料质量及后续部署的推理摊销成本。

面向用户交互的大模型还需与人类意图对齐。InstructGPT 结合指令监督微调（SFT）与基于人类反馈的强化学习（RLHF），塑造模型的指令遵循与价值对齐表现\cite{ouyang2022instructgpt}。后训练阶段引导模型在保证事实性的前提下，输出安全、结构规范且有帮助的回答。

\subsection{语言模型的能力扩展}
语言模型的能力边界持续外延：视觉指令微调将跨模态视觉特征投射至语言表示空间\cite{liu2023visualinstruction}；ReAct 范式将思维推理与外部工具调用交替交织\cite{yao2023react}；DeepSeek-R1 等工作表明强化学习可激发长思维链的自主探索与可验证推理能力\cite{deepseek2025r1}。

针对高难度数学与逻辑问题，系统可通过测试时计算扩展（Test-time Compute Scaling）为模型分配生成、蒙特卡洛树搜索或外置代码解释器验证资源。模型生成动作由外围安全运行时隔离执行，形成自主智能体架构。

\begin{table}[htbp]
\centering
\caption{语言模型建模范式演进与特征对比}\label{tab:intro-paradigms}
\begin{tabularx}{\linewidth}{@{}p{22mm}p{34mm}p{20mm}X@{}}
\toprule
建模范式 & 上下文与记忆机制 & 训练并行度 & 核心特征与代表模型 \\\midrule
$n$-gram & 局部滑动窗口（$n-1$ 词） & 无需参数训练 & 容易遭遇数据稀疏；无表示泛化。 \\
RNN / LSTM & 递归隐状态 $\vect{h}_t=f(\vect{h}_{t-1},x_t)$ & 串行依赖 & 摆脱固定窗口；但长程梯度衰减，无法并行。 \\
自编码（BERT） & 双向全互联自注意力 & 序列完全并行 & 适合文本理解与抽取；不适于长文本生成。 \\
自回归（GPT） & 因果单向自注意力 & 训练完全并行 & 概率链式分解严格自洽；天然适合开放生成。 \\
\bottomrule
\end{tabularx}
\end{table}

\subsection{序列模型的局限性}
模型利用序列信息的范围和效果取决于可见位置、状态容量及信息传播路径。固定窗口模型只读取预先规定的邻近位置，窗口之外的信息不能进入当前预测；扩大窗口又会增加输入维度和需要覆盖的上下文组合。RNN 用递归状态摆脱了固定窗口，却需要把不断增长的历史压缩到固定维度状态中。信息和梯度从早期位置传到远处位置时经过的变换随距离增加，因而容易出现记忆衰减、梯度消失或梯度爆炸。LSTM 的门控与加性记忆路径缓解了这些问题，但没有消除状态容量和长程训练信号的限制。

递归结构还形成沿时间步的执行依赖：计算第 $t$ 个状态之前必须先得到第 $t-1$ 个状态，因此长序列训练难以充分并行。自注意力把任意两个可见位置之间的信息路径缩短为直接交互，并允许训练时同时计算多个位置；密集自注意力需要构造位置两两交互，其计算量和中间结果规模随序列长度呈二次增长。扩大上下文长度因此同时增加计算、显存和数据需求，相关信息的提取效果还受位置分布、干扰内容及长距离训练信号影响。

因果语言模型还有训练与生成之间的并行性差异。训练时已知完整目标序列，可以在因果掩码下\emph{并行计算}各位置的条件分布；生成时，第 $t$ 个词元的实际输入包含前面已经选出的词元，必须\emph{逐步自回归决策}。键值缓存能够避免重复计算历史位置，却不能消除这种\emph{决策依赖}，缓存本身也随上下文增长并占用显存与访存带宽。序列建模的系统权衡需从注意力可见域、状态表征容量、梯度传播路径、时序执行依赖与显存访存开销等多维约束综合评估。\bookxref{ch:tokenization}进一步解释固定窗口、RNN 与 LSTM，\bookxref{ch:attention}推导注意力的计算和存储代价，\bookxref{ch:transformer}与\bookxref{ch:incremental-decoding-cache}分别讨论因果训练和增量生成。

\section{规模定律与涌现能力}

\label{sec:intro-scaling-emergence}
\subsection{规模定律的适用范围}
\term{规模定律}{Scaling Law}{}是对模型、数据、计算与测量结果之间关系的经验描述。对语言模型，常见研究对象是在可比数据分布上测得的平均验证交叉熵。一个用于理解参数与数据约束的拟合形式为
\begin{equation}
 L(N,D)\approx E+\frac{A}{N^{\alpha}}+\frac{B_s}{D^{\beta}},
 \qquad A,B_s,\alpha,\beta>0.
\end{equation}
其中 $N$ 为参数量，$D$ 为训练词元数，$E$ 为拟合底座，其余系数由实验估计。在拟合适用范围内，扩大参数或数据可以降低相应的损失余量；只扩大其中一项，另一项仍可能限制收益。$E$ 是依据实验数据估计的渐近项，其数值依赖数据分布与拟合模型。

Kaplan 等人报告了特定实验条件下损失随规模变化的幂律规律；Chinchilla 工作进一步研究固定训练计算下的参数与数据配置，得到其计算最优范围内两者应近似同比例增长的结论\cite{kaplan2020scaling,hoffmann2022chinchilla}。这些结果依赖各自的训练设置与拟合范围；把它们当作适用于所有模型的固定词元参数比会超出适用范围。模型结构、语料质量、重复使用数据及推理服务成本变化，都可能改变预算选择。

对常见密集模型的粗略训练预算，可写成 $C\approx\kappa ND$，其中 $C$ 是训练计算量，$\kappa$ 汇总结构与计算口径相关的系数。在固定 $C$ 下增加 $N$ 就要压缩 $D$，更大模型能否带来更低损失取决于两者的配比。\bookxref{ch:pretraining-scaling}进一步推导拟合、计算最优分配和外推验证；本章先建立一条判断：\keyconcept{规模定律用于有条件的预测与规划，它的被解释变量是损失，任务正确率由具体任务指标另行测定。}

\begin{note}[规模定律的经验性质与任务外推边界]
规模定律建立在特定架构、分布和测量指标（通常为验证交叉熵）之上。损失平滑下降并不必然等同于特定下游任务正确率呈线性改善；低交叉熵是良好任务表现的必要倾向，而非绝对保证。工程规划时需将损失预测与下游评测基线结合验证。
\end{note}

\subsection{涌现能力的度量}
\term{涌现能力}{Emergent Abilities}{}在 Wei 等人的2022年研究中，指一些在较小模型中未观察到、在较大模型中出现，且难以从小规模表现外推的任务能力\cite{wei2022emergent}。这一说法描述给定模型系列、任务与测量方式下的现象。“小模型未观察到”由测试集大小、随机基线、提示方式和评价粒度共同决定，内部表示的变化则通过表征分析考察。

2023年 Schaeffer 等人指出，一些表面上的突跃可以由非线性或不连续的评价指标产生；改变度量方式后，部分能力曲线呈现平滑改善\cite{schaeffer2023mirage}。分析能力曲线时，任务性质与度量方式都要纳入考察。任务要求完整答案正确时，精确匹配有实际意义；判断内部机制是否突变，则要依靠连续度量和表征分析。

例如，设答案包含10个位置，每个位置独立且以相同概率 $q$ 正确，则整串完全正确的概率为 $q^{10}$。当 $q$ 从 $0.5$ 增至 $0.8$ 时，逐位置正确率从50\%升至80\%，整串正确率却从约0.10\%升至10.74\%。少量测试题上，前者可能长期测为零，随后才出现明显非零分数。位置间错误相关时，整串正确率还取决于它们的联合分布。

研究涌现时，应同时报告严格任务指标与更细粒度的诊断指标，并在疑似转折附近增加规模点、重复实验和不确定性估计。模型之间的数据、训练目标、提示或解码方式若同时变化，观察到的差异就不能全部归因于规模。数据污染与样本记忆属于需要单独排查的来源。

\subsection{平均损失和任务能力的关系}
平均损失汇总大量预测位置，任务评分则可能只关心某个领域，或要求多个步骤同时正确。前者平滑下降时，后者仍可能因阈值和组合要求而快速变化，平均改善也会掩盖少数任务的退步。规模定律与涌现研究由此关注不同层次的观测。

\emph{训练阶段计算}与推理阶段\emph{测试时计算}具有不同的扩展规律。增加训练词元通过梯度更新提升基座模型的泛化能力；推理期分配额外的候选采样、树搜索或链式验证预算，则是对固定参数模型的计算增强，其边际收益取决于解空间剪枝与验证信号的准确度。

\section{语言模型的概率建模}

模型通常不能直接获知真实语言分布，只能观察有限语料。设数据集为 $D=\{x^{(1)},\ldots,x^{(n)}\}$，一种基本学习方式是选择使观测数据似然尽可能大的参数，即最大似然估计：
\begin{equation}
 \theta^*\in\argmin_\theta\left[-\sum_{i=1}^{n}\log p_\theta(x^{(i)})\right].
\end{equation}
将序列概率分解代入，便得到逐词元负对数概率的求和。将其表述为可微目标后，优化算法可据此计算梯度并更新参数。

最大似然估计促使模型逼近训练语料的经验分布；语料中的事实谬误、逻辑缺陷与偏见同样会被模型拟合。生成\emph{流利文本}并不等同于保证\emph{事实正确}，事实真实性需通过知识检索、逻辑核验及外围控制机制协同保障。

自监督学习从数据自身构造目标，下一词元预测就是其中一种方式。它省去了人工逐词元标注，数据的采集、筛选、混合与组织仍然决定模型能力的边界。

\begin{example}[局部决策与序列决策]
得到概率后，还需要决定如何选择输出。考虑只生成两个词元的简化任务。首步只有 $a$、$b$ 两个候选，其概率分别为 $0.6$ 和 $0.4$。第二步候选为 $u$、$v$，条件概率如下。
\begin{table}[htbp]
\centering
\caption{两步生成条件概率分布示例}
\label{tab:intro-two-step-example}
\begin{tabular}{cccc}
\toprule 首词元 & 首步概率 & $p(u\mid\text{首词元})$ & $p(v\mid\text{首词元})$\\\midrule
$a$ & $0.6$ & $0.5$ & $0.5$\\
$b$ & $0.4$ & $0.9$ & $0.1$\\\bottomrule
\end{tabular}
\end{table}
由乘法规则，四条完整序列的概率为
\begin{align*}
 p(a,u)&=0.6\times0.5=0.30,&p(a,v)&=0.30,\\
 p(b,u)&=0.4\times0.9=0.36,&p(b,v)&=0.04.
\end{align*}
四者之和为一。逐步贪心策略先选择局部概率最大的 $a$，后续无论如何选择均只能得到 $0.30$ 的序列概率；而全局联合概率最大的路径为 $(b,u)$，概率达到 $0.36$。该例表明逐步贪心搜索无法保证\emph{全局最优}，全局求解需依赖束搜索（Beam Search）等序列解码算法。

生成概率刻画的是模型对语料统计模式的拟合，答案正确性则需由具体任务目标检验。获取高概率序列之后，仍需验证其是否满足逻辑与事实要求。
\end{example}

\section{语言模型系统的组成}

\subsection{模型结构}
网络结构规定输入信息如何映射为输出分数，包括表示维数、模块拓扑、注意力掩码与非线性变换。训练算法规定批次采样、损失计算、反向传播与优化器更新规则。相同网络在不同语料与优化目标下将获得不同能力；不同网络在相同目标下亦具有独特的计算复杂度与归纳偏置。

Transformer 借助自注意力矩阵交换序列位置信息，再经由前馈子层实施逐位置非线性变换。编码器、因果解码器和编码器—解码器均由这些组件堆叠而成，各架构具有不同的信息可见性与监督目标。

模型规模由多项互补指标共同刻画：总参数量表征模型的函数容量与持久化存储需求；激活参数量（如混合专家网络中单次前向激活的子集）决定单步浮点运算量；上下文窗口界定注意力机制可直接感知的序列跨度；训练词元总数则决定数据多样性与训练阶段的总计算预算。

\begin{table}[htbp]
\centering
\caption{不同神经网络结构的归纳偏置与特性对照}\label{tab:intro-structures}
\begin{tabularx}{\linewidth}{@{}p{32mm}XX@{}}
\toprule 结构 & 核心机制 & 适用性与限制\\\midrule
线性模型／树模型 & 线性打分／递归划分特征空间 & 可形成表格任务基线；能力取决于特征与划分。\\
多层感知机（MLP） & 全连接映射与逐点非线性组合 & 适合固定维度表示，不自动利用空间邻接。\\
卷积神经网络（CNN） & 局部连接和空间共享参数 & 利用局部结构；跨远距离关系需多层或其他机制。\\
循环神经网络（RNN） & $\vect{h}_t=f_\theta(\vect{h}_{t-1},x_t)$ 的递归状态 & 顺序更新；历史压缩在状态中，路径长度随时间增长。\\
Transformer & 内容相关注意力与逐位置变换 & 可按可见性构造不同序列模型，计算代价随连接数增长。\\
扩散模型（Diffusion） & 学习逐步去噪或对应连续动力学 & 在数据或潜空间生成；采样通常需多步计算。\\
\bottomrule
\end{tabularx}
\end{table}
长短期记忆通过输入门、遗忘门与输出门控制隐状态流动，递归时序依赖依然存在。局部卷积、递归隐状态与自注意力分别引入不同的\term{归纳偏置}{Inductive Bias}{}，即在有限数据约束下对特定函数形式的先验偏好。

\subsection{学习范式}
\label{sec:intro-learning-map}
学习范式的划分依据在于监督信号的来源与反馈机制：预训练阶段采用自监督方式，根据文本前缀自回归预测后续词元；指令微调阶段引入人工标注或合成示范数据进行\term{监督学习}{Supervised Learning}{}；对齐阶段则借助人类或模型评判构建偏好奖励，利用\term{强化学习}{Reinforcement Learning}{RL}优化策略；\term{无监督学习}{Unsupervised Learning}{}则侧重刻画分布隐变量或聚类几何结构。

\subsection{模型系统的生命周期}
训练系统将语料转化为参数权重，应用系统将用户请求转化为生成响应。两者共享模型制品，但维护不同的执行状态与容灾要求：训练系统聚焦\emph{吞吐与一致性}（数据加载、梯度同步、检查点保存），应用系统聚焦\emph{延迟与可靠性}（排队时延、首字时间、工具幂等、输出安全）。
\begin{figure}[htbp]
\centering
\includegraphics[width=\linewidth]{\subfix{../../figures/theory/revision-learning-serving.pdf}}
\caption[训练与应用的数据路径]{训练与应用的数据路径。实线表示数据或制品流动，虚线表示评估反馈；反馈进入下一轮修订，不自动改变正在服务的模型。}\label{fig:lifecycle}
\end{figure}
图\ref{fig:lifecycle}中的模型制品不仅包含权重张量，还包含分词器配置、模型超参数架构与输入模板协议。外围适配层或多模态连接器需与基座协同加载。

大模型训练系统由加速器算力与显存、主机数据预处理与调度、高速互联网络以及分布式存储构成端到端流水线。任一环节吞吐受限都会导致加速核心闲置，因此集群算力规划必须基于端到端通信与 I/O 拓扑设计，而非仅依赖单卡理论峰值算力。

\subsection{模型的外部能力扩展}
检索增强生成（RAG）通过检索外部知识库获取上下文片段，作为模型生成的显式输入。工具调用使语言模型能够触发查询 API、数值计算或数据库事务。前者拓展了模型的即时上下文窗口，后者拓展了系统的外部执行能力。智能体进一步根据环境观测动态编排动作序列，涉及目标规划、短期记忆维护与异常恢复。

工具调用由安全执行器校验参数合法性与执行权限；生成的引用信息需与原始检索文档严格对齐，以保证论断可溯源。

\section{学习路径}

共同基础围绕一次模型请求建立概念链：文本如何编码为词元张量，自注意力如何计算上下文条件分布，解码策略如何搜寻与终止输出，以及性能与安全性如何量化评测。

\begin{figure}[htbp]
\centering
\includegraphics[width=\linewidth]{\subfix{../../figures/theory/intro-technology-stack.pdf}}
\caption[大语言模型技术层次体系]{大语言模型技术层次体系。涵盖算力基础设施、模型架构与表示、训练与对齐、推理服务及应用智能体五个核心层次。}\label{fig:intro-technology-stack}
\end{figure}

应用智能体系统聚焦任务达成与外部交互，经由上下文学习、向量检索与 RAG 进入领域落地，并深入智能体运行时架构。模型训练方向聚焦数据驱动下的参数学习，由监督微调、参数高效微调延伸至预训练、强化学习与可验证推理。

训练与推理系统聚焦底层资源调度与计算优化，涵盖硬件通信拓扑、显存预算分析、分布式并行、增量解码缓存管理与服务批处理调度。

全书各方向通过模型制品、通信协议与评测证据相互衔接：应用系统依托固定权重构建外部运行时，适配技术聚焦特定任务下的参数演化，系统工程则围绕计算与显存效率展开底层优化。

语言模型的理论推导与工程实现建立在经典统计学习与现代深度学习的基础之上。在进入大模型特有的自注意力机制、预训练与微调、分布式并行及智能体系统之前，若需系统巩固数理推导（如概率图模型、极大似然估计与优化目标），可参阅李航《统计学习方法》\cite{LiTongJiXueXiFangFaDi2Ban2019}；若需建立张量计算、自动微分与深度网络基础模块的代码实践与直觉，可参阅张阿斯顿、李沐等《动手学深度学习》\cite{zhangDiveDeepLearning2022}。具体章节路线与先修对应说明见书前《阅读说明》。

\begin{table}[htbp]
\centering
\caption{核心数学符号与张量记号规范}\label{tab:intro-notations}
\begin{tabularx}{\linewidth}{@{}l p{30mm} X@{}}
\toprule
记号符号 & 类型／维度 & 含义与全书统一规范约定 \\\midrule
$x, y, t$ & 标量 & 离散词元编号、步长、时钟索引或温度超参数 \\
$\vect{x}, \vect{h}, \vect{z}$ & 向量（粗斜体） & 隐层激活状态、词嵌入向量或模型预测对数概率 Logits \\
$\mat{W}, \mat{E}, \mat{A}$ & 矩阵（大写粗体） & 线性变换权重、词表嵌入矩阵或低秩适配矩阵 \\
$B, S, D$ & 正整数标量 & 批量大小（Batch Size）、序列长度（Sequence Length）、隐藏维度（Hidden Dim） \\
$H, d_k, d_v$ & 正整数标量 & 注意力头数（Heads）、键维度（Key Dim）、值维度（Value Dim） \\
$V, |\mathcal{V}|$ & 正整数标量 & 词表大小（Vocabulary Size） \\
$p_\theta(x_t \mid x_{<t})$ & 条件概率分布 & 参数为 $\theta$ 的模型在当前前缀已知时对下一词元的预测概率 \\
\bottomrule
\end{tabularx}
\end{table}

\chapterreferences
\myendchapterdocument
